\section*{Chapitre \ref{ch:m2:loans-clos-and-securitisation} --- Prêts, CLO et titrisation}

\begin{solution}{exo:m2:loans-clos-and-securitisation:1}
$500/365 = 1.370$, $500/395 = 1.266$, $500/425 = 1.176$ et $500/450 = 1.111$.
\end{solution}

\begin{solution}{exo:m2:loans-clos-and-securitisation:2}
Le prêt verse un taux variable plus un spread fixe, réajusté à chaque période à partir d'un taux court qui suit le taux directeur
de la banque centrale~; le coupon de l'obligation a été fixé à l'émission.
\end{solution}

\begin{solution}{exo:m2:loans-clos-and-securitisation:3}
Les covenants de maintien~: le droit d'agir, de renégocier ou d'exiger le remboursement immédiat, lorsque les
ratios de l'emprunteur se dégradent, avant qu'il ne manque un paiement. Les prêteurs apprennent
les difficultés plus tard et ont moins voix au chapitre jusqu'au défaut.
\end{solution}

\begin{solution}{exo:m2:loans-clos-and-securitisation:4}
Le pair peut baisser de $500 - 1.18 \times 395 = 33.9$~millions de dollars~; avec une perte de 30~\% par
défaut, cela fait 113~millions de dollars de défauts, soit 22.6~\% du portefeuille.
\end{solution}

\begin{solution}{exo:m2:loans-clos-and-securitisation:5}
La règle s'applique aux titriseurs, qui transfèrent des actifs dans une \omterm{def:m2:loans-clos-and-securitisation:sec}{titrisation}~; un
gestionnaire de CLO en marché ouvert n'octroie ni ne détient les prêts avant que l'\omterm{def:m2:loans-clos-and-securitisation:sec}{entité ad hoc} ne les achète, il n'a donc rien à retenir en ce sens.
La rétention des risques vise le modèle «~octroyer pour distribuer~»~: un prêteur qui revend
tout ce qu'il octroie a peu de raisons de
prêter avec soin.
\end{solution}

\begin{solution}{exo:m2:loans-clos-and-securitisation:6}
Les prêts rapportent 7.5~\% sur 500~millions de dollars~; les billets coûtent en moyenne 5.81~\% sur 450
millions de dollars, et les frais 0.45~\% sur le tout. La \omterm{def:m2:loans-clos-and-securitisation:tests}{tranche equity}, 50~millions de dollars, touche
l'écart sur l'ensemble du portefeuille~: $37.5 - 2.25 - 26.14 = 9.11$~millions de dollars, soit 18.2~\% de
son investissement, car sa propre mise, faible, rapporte le spread gagné sur un montant dix fois plus élevé.
\end{solution}

\begin{solution}{exo:m2:loans-clos-and-securitisation:7}
5.35~\% par an au lieu de 4.78~\%. Après la période de réinvestissement, les recouvrements remboursent les
billets AAA au lieu d'acheter de nouveaux prêts, ce qui fait baisser les billets plus vite que le pair et
fait monter les ratios de couverture, si bien que les tests échouent plus tard.
\end{solution}

\begin{solution}{exo:m2:loans-clos-and-securitisation:8}
Les billets AAA du CLO sont notés sur la structure~: leur sécurité vient de la subordination
et de tests portant sur un portefeuille de prêts spéculatifs~; elle dépend de la façon dont les défauts
du portefeuille sont corrélés, ce qu'une obligation d'une seule entreprise ne fait pas. Leur prix peut
chuter bien davantage lors d'une crise de crédit, quand corrélation et dégradations de notation augmentent ensemble, et ils
sont moins liquides.
\end{solution}

\begin{solution}{pb:m2:loans-clos-and-securitisation:1}
\textbf{1.} 37.5~millions de dollars~; les frais 2.25~millions de dollars, les billets 26.14~millions de dollars.
\textbf{2.} 9.11~millions de dollars, 18.2~\%.
\textbf{3.} 19.5~\% par an.
\textbf{4.} 27.5~millions de dollars ($500 - 1.05 \times 450$).
\textbf{5.} 91.7~millions de dollars, 18.3~\% du portefeuille.
\textbf{6.} 13.1~\%, 5.0~\% et $-5.8\%$.
\textbf{7.} 2.0, 12.5 et 28.4 millions de dollars.
\textbf{8.} 4.78~\% par an.
\textbf{9.} Au trimestre 19, lorsque le pair du portefeuille a baissé de 33.6~millions de dollars.
\textbf{10.} 6.5~millions de dollars sur ses 25~millions de dollars ne sont pas remboursés, et la \omterm{def:m2:loans-clos-and-securitisation:tests}{tranche equity} ne reçoit
rien à partir du trimestre 8.
\textbf{11.} Il porte le seuil à 5.35~\%~: les recouvrements remboursent alors les billets.
\textbf{12.} Chaque défaillance transforme les flux subordonnés en remboursements des billets AAA, ce qui
fait baisser leur encours plus vite que le portefeuille ne perd de pair.
\textbf{13.} Chaque défaut fait baisser le pair de la perte, 30~\% pour le prêt contre 60~\% pour
l'obligation, il consomme deux fois moins de coussin.
\textbf{14.} Vendre les prêts fragiles avant leur défaut, ou acheter des prêts sous le pair, que les
tests peuvent compter au pair~; la première cristallise des pertes, la seconde ajoute du risque, et
les deux sont limitées par la documentation.
\textbf{15.} Ce mécanisme restitue la part du principal revenant à la \omterm{def:m2:loans-clos-and-securitisation:tests}{tranche equity}~; si les prêts doivent être vendus
sous le pair, comme sur un marché sous tension, ce dernier versement se réduit ou disparaît.
\textbf{16.} Avec les taux de défaut réalisés sur les \omterm{def:m2:loans-clos-and-securitisation:loan}{prêts à effet de levier}, y compris les
échanges de dette en détresse, lors des récessions passées~; les défauts se concentrent pendant une
récession, si bien que quelques années nettement au-dessus de la moyenne causent plus de dégâts que le même total réparti uniformément.
\textbf{17.} La \omterm{def:m2:loans-clos-and-securitisation:tests}{tranche equity}~: des investisseurs capables de supporter de fortes pertes et qui valorisent
l'excédent de spread~; les billets AAA~: des investisseurs qui ont besoin de fortes notations et d'actifs à taux variable, tels
que les banques et les assureurs.
\textbf{18.} Acheter une protection sur un \omterm{def:m2:credit-indices-and-tranches:index}{indice de crédit} high yield ou sur sa tranche junior, qui
perdent dans la même vague de défauts~; la couverture est imparfaite, car les prêts du CLO ne sont pas
les émetteurs de l'indice.
\textbf{19.} Résultat nommé~: \emph{le test qui se déclenche} coupe la \omterm{def:m2:loans-clos-and-securitisation:tests}{tranche equity} à un
taux de défaut constant de 4.78~\% par an, dès le trimestre 19, lorsque le pair a baissé
de 33.6~millions de dollars.
\textbf{20.} Il prélève des liquidités sur la \omterm{def:m2:loans-clos-and-securitisation:tests}{tranche equity} et les verse aux billets seniors dès que
les pertes entament le coussin.
\end{solution}

\begin{solution}{iq:m2:loans-clos-and-securitisation:1}
Une \omterm{def:m2:loans-clos-and-securitisation:sec}{entité ad hoc} qui détient un portefeuille de \omterm{def:m2:loans-clos-and-securitisation:loan}{prêts à effet de levier} géré activement, financée par des billets notés
de AAA à BB et par une \omterm{def:m2:loans-clos-and-securitisation:tests}{tranche equity} non notée~; les flux de trésorerie paient les billets dans l'ordre et le
reste à la \omterm{def:m2:loans-clos-and-securitisation:tests}{tranche equity}. Les banques, assureurs, fonds et caisses de retraite détiennent les billets notés~; la
\omterm{def:m2:loans-clos-and-securitisation:tests}{tranche equity} revient aux investisseurs capables de supporter les premières pertes.
\iqlookfor{SPV, gestionnaire, tranches, investisseurs selon le risque.}
\end{solution}

\begin{solution}{iq:m2:loans-clos-and-securitisation:2}
À chaque date de paiement, le ratio entre le pair des prêts et les billets jusqu'à une classe donnée, et celui
des intérêts encaissés sur les intérêts dus, sont comparés à des seuils de déclenchement~; si l'un échoue, la
trésorerie qui serait allée plus bas rembourse les billets seniors jusqu'au rétablissement du ratio.
Les défauts, les pertes et la dégradation de notation des prêts font bouger les ratios.
\iqlookfor{ratios de pair et d'intérêts, seuils, détournement des flux.}
\end{solution}

\begin{solution}{iq:m2:loans-clos-and-securitisation:3}
Modéliser les défauts, les recouvrements, les remboursements anticipés et les taux du portefeuille, par scénarios ou
simulation, avec le détail prêt par prêt si disponible~; faire passer la \omterm{def:m2:loans-clos-and-securitisation:clo}{cascade des paiements} sur chaque trajectoire~; la
valeur de la \omterm{def:m2:loans-clos-and-securitisation:tests}{tranche equity} est la trésorerie résiduelle actualisée, à un taux reflétant son risque, et
sa sensibilité aux hypothèses de défaut et de recouvrement constitue le principal risque. Il faut vérifier le
résultat aux prix du marché secondaire.
\iqlookfor{scénarios passés par la cascade réelle.}
\end{solution}

\begin{solution}{iq:m2:loans-clos-and-securitisation:4}
Une tranche d'indice est synthétique, statique, d'échéance fixe et valorisée chaque jour à partir des
spreads de l'indice~; une tranche de CLO est une position en liquidités, adossée à un portefeuille de prêts géré et évolutif, avec
réinvestissement et des tests qui déplacent les liquidités entre classes, et se traite rarement. Son risque
dépend du gestionnaire et de la documentation, et non seulement de la corrélation des défauts.
\iqlookfor{statique contre géré, tests, liquidité.}
\end{solution}

\begin{solution}{iq:m2:loans-clos-and-securitisation:5}
La hausse de l'effet de levier et l'affaiblissement des covenants chez les emprunteurs, la dette à taux variable qui les rend
vulnérables à la hausse des taux, concentration dans certains secteurs, et la chaîne des
détenteurs (CLO, fonds, assureurs et banques) par lesquels les pertes se propagent~; les fonds qui
promettent une liquidité sur des prêts illiquides peuvent faire face à des ruées.
\iqlookfor{fragilité des emprunteurs et lieu où se concentrent les pertes.}
\end{solution}

\begin{solution}{iq:m2:loans-clos-and-securitisation:6}
Encoder la cascade de chaque opération comme des données (classes, tests, seuils, frais, dates) exécutée par un
moteur unique plutôt que du code propre à chaque opération~; vectoriser sur les scénarios~; partager les flux
de trésorerie au niveau des prêts entre opérations qui détiennent les mêmes prêts~; exécuter en parallèle les opérations et les
nœuds de calcul~; valider chaque opération par rapport aux rapports du trustee.
\iqlookfor{cascade en tant que données, réutilisation, parallélisme, validation.}
\end{solution}
